\documentclass[12pt,a4paper]{article}

\usepackage[T1]{fontenc}
\usepackage[utf8]{inputenc}
\usepackage[margin=2.5cm]{geometry}
\usepackage{setspace}
\usepackage{csquotes}
\usepackage{hyperref}

\begin{document}
	
	\title{What world do they want?\\
		\large The ideological structure of the contemporary Western Left}
	\date{June 2026}
	\maketitle
	
	\begin{abstract}
		\noindent
		Contemporary left-wing political movements in Europe and the United States exhibit
		a pattern of contradictions that resist explanation through the conceptual vocabulary
		of twentieth-century socialism: stated commitments to democracy, secularism, and
		gender equality coexist with observable alliances with forces antithetical to each of
		those commitments. This article attempts a structural account of that pattern. Three
		complementary explanatory frameworks are proposed, identitarian coalitionism, the
		post-Marxist substitution of the revolutionary subject, and the psychosocial function
		of collective guilt, and integrated into a synthetic description of the positive
		project, such as it is, that underlies contemporary left-wing activism. The analysis
		concludes that the observable incoherence is not accidental but follows from a
		coherent internal logic, and that the most serious risk it generates is not for its
		declared adversaries but for the movement itself.
		
		\medskip
		\noindent\textbf{Keywords:} identity politics, intersectionality, post-Marxism,
		Western left, political coalitions, cultural relativism, collective guilt, equity.
	\end{abstract}
	
	\section{The problem of apparent incoherence}
	
	For an observer who came of age politically in the twentieth century, the behaviour
	of a substantial portion of today's Western left presents a genuine puzzle. The
	same political actors who invoke democracy align themselves with movements that
	are openly hostile to democratic procedure. Those who assert the primacy of the
	secular state extend systematic indulgence to a religious fundamentalism they would
	condemn without hesitation if it were Christian. Those who most loudly advocate
	gender equality adopt a studied silence, or, worse, an active apologetics, in the
	face of practices within certain immigrant or postcolonial communities that they
	would characterise as patriarchal oppression in any other context.
	
	The temptation is to attribute this pattern to hypocrisy, to tactical opportunism,
	or simply to a collective failure of reasoning. That temptation should be resisted.
	The pattern is too consistent and too widespread to be explained by individual
	bad faith. A more productive approach is to ask whether there is an internal logic
	that makes these positions coherent from within the framework that generates them,
	even if that framework itself rests on contestable premises.
	
	This article argues that there is such a logic, and that it derives from three
	interlocking structural features of contemporary left-wing thought: a shift in the
	primary axis of political analysis from class to identity; a relativism about
	cultural values that is asymmetrically applied; and a psychosocial dynamic in which
	collective guilt functions as the constitutive identity of a political community.
	Together these features produce what might be called a negative utopia, a political
	project defined far more precisely by what it opposes than by what it would
	construct.
	
	It should be noted at the outset that the phenomenon described here does not
	encompass the whole of the left. There exist left-leaning individuals and traditions
	who identify these contradictions clearly, denounce them, and are for that reason
	marginalised or accused of having defected to the right. The analysis that follows
	pertains to the identitarian and postcolonial current that has achieved institutional
	hegemony in Western universities, major media organisations, and large NGOs over the
	past two to three decades.
	
	\section{From Class to Identity}
	
	The historical left, whether social-democratic or Marxist, organised its worldview
	around an economic axis. Society was divided between those who owned the means of
	production and those who sold their labour; political conflict was ultimately a
	conflict over the distribution of the surplus generated by that relationship. The
	primary categories were class categories, and the primary subject of emancipation
	was the working class, defined by its position in the productive process regardless
	of ethnicity, religion, or culture.
	
	The contemporary left has replaced this axis with a different one. Shaped by
	post-structuralist theory from the 1970s onward, and more recently by the framework
	of intersectionality, it reads social reality through the lens of power and
	identity. The fundamental division is no longer between bourgeoisie and proletariat
	but between those who hold systemic power and those who are subject to it. On the
	side of power: Western civilisation, whiteness, Christianity, heterosexuality,
	capitalism, and the male sex. On the side of subjection: ethnic minorities, the
	global South, Islam in Western contexts, LGBTQ+ communities, and women, understood
	not as individuals but as members of identity groups.
	
	The crucial feature of this framework is that the axis of oppression is determined
	structurally, not by the content of the values held by any given group. A religious
	fundamentalist who belongs to a minority community in a Western country or to a
	formerly colonised nation is categorised, first and foremost, as occupying a
	position of structural subjection relative to Western hegemony. The values he holds
	regarding women, homosexuals, or apostates are secondary to his structural position.
	The coalition is organised not around shared positive values but around a shared
	structural relationship to the dominant order.
	
	This is why the contradictions described in the introduction are not, from within
	the framework, contradictions at all. The ordering criterion is position in the
	power hierarchy, not content of belief. Anyone who is an Other relative to the
	Western mainstream is a potential ally; the specific content of that alterity is
	not the primary variable. The apparent incoherence resolves, within the framework,
	into a rigorous if contestable consistency.
	
	\section{The architecture of guilt}
	
	A second structural feature concerns the treatment of historical responsibility.
	For the contemporary left, the colonial history of the West is not a concluded
	chapter that can be weighed against other chapters in a balanced historical account.
	It is, instead, a permanent structure: the current wealth, stability, and global
	influence of Western nations are understood as having been extracted from the
	colonised world and as being actively maintained at that world's expense. Historical
	injustice does not recede; it reproduces itself in the present through the
	institutions, cultural norms, and relative prosperity of Western societies.
	
	This framing has a direct consequence for the application of normative standards
	across cultural boundaries. To hold non-Western communities to Western normative
	standards, to criticise patriarchal practices, theocratic governance, or the
	treatment of religious minorities in majority-Muslim societies, is, in this
	framework, an act of cultural imperialism: the imposition of norms that are
	themselves products of a hegemonic and guilty civilisation. The same standards
	that apply universally within the West cannot be applied universally across cultural
	boundaries without reproducing the logic of colonial domination.
	
	The result is an asymmetric application of principles that appears, from the outside,
	as hypocrisy. Sexism is condemned without qualification when it occurs within Western
	societies, but contextualised, explained, and ultimately excused when it occurs
	within communities framed as postcolonial. This is not inconsistency; it is the
	application of a coherent, if disputable, normative hierarchy in which
	anti-colonialism ranks above secularism and anti-racism ranks above gender equality.
	
	The guilt is, as a consequence, structurally interminable. It cannot be discharged
	by any finite act of reparation, because it is not located in discrete historical
	events but in the ongoing structure of Western power. The only conceivable terminus
	is the dismantling of that structure, which is, precisely, the political objective.
	
	\section{Collective guilt as identity}
	
	The third structural feature operates at the psychosocial rather than the
	theoretical level. Pascal Bruckner, in \emph{The Tyranny of Guilt}, and subsequent
	critics of what has come to be called woke ideology, have identified a dynamic in
	which the acknowledgement of collective guilt functions not merely as a moral
	position but as the constitutive identity of a political community.
	
	The mechanism is as follows. The historical indictment of Western civilisation, for
	colonialism, slavery, patriarchy, and the various other crimes attributed to it, is
	not experienced by adherents primarily as an external factual claim subject to
	evidence and revision. It is experienced as a marker of moral seriousness, of
	intellectual sophistication, and of belonging to a community defined by its
	rejection of complacency about Western privilege. To question the indictment, or to
	apply it with nuance, or to note cases that do not fit the framework, is experienced
	not as a contribution to accuracy but as a betrayal of the community and a
	regression to comfortable ignorance.
	
	This dynamic explains what might otherwise appear as an irrational impermeability
	to counterevidence. Evidence is not irrelevant because adherents are irrational; it
	is irrelevant because the primary function of the framework is not predictive or
	explanatory but identity-constitutive. Colonised countries in Asia that achieved
	rapid economic development without the pathologies attributed to Western colonialism
	do not falsify the framework; they are simply not assimilated into it, because the
	framework is not, at its core, an empirical hypothesis about the causes of
	underdevelopment. It is a story about who one is and who one opposes.
	
	The guilt, in this reading, cannot be expiated because expiation would dissolve the
	identity. A community that had genuinely discharged its historical debt would lose
	the central organising principle of its collective self-understanding. The
	interminability of the guilt is thus not incidental but functional.
	
	\section{The positive project: a radically egalitarian Utopia}
	
	Having described the internal logic of the framework, one can now address the
	central question more directly: what positive world, if any, do the holders of this
	framework want?
	
	The aspiration, when it is made explicit, is a radically egalitarian utopia: a
	social order in which every hierarchy of power has been dismantled. This means a
	world without significant economic inequality, without the privileged position of
	any cultural tradition, without the nation-state as the primary unit of political
	organisation, and without the normative weight currently attached to the family
	structures and religious traditions of the Western mainstream. It is a world in
	which the historical debts of Western civilisation have been expiated through a
	comprehensive redistribution of power, wealth, and cultural authority.
	
	The connection between this positive vision and the coalitional and guilt-based
	structures described above is direct. The existing foundations of Western
	society, capitalism, the Judeo-Christian cultural inheritance, the nation-state,
	the bourgeois family, are held to be intrinsically and irremediably corrupt. They
	cannot be reformed; they must be deconstructed. Any force that weakens or
	destabilises them is therefore, at least provisionally, an ally in the project of
	deconstruction, regardless of whether that force shares the positive vision of the
	egalitarian utopia that is supposed to follow.
	
	This is the structural role played by alliances that appear, from the outside, to be
	self-defeating. Radical Islamism, in this framework, is not a value-based ally; it
	is a tactical instrument, consciously or unconsciously, whose value consists in
	its antagonism toward the existing order. It is an updated version of what Lenin
	called the fellow traveller: useful for the demolition phase, to be managed in the
	construction phase that comes after.
	
	The obvious difficulty, which the historical record underlines with considerable
	force, is that the construction phase has a way of not arriving on schedule. The
	Iranian revolution of 1979 offers the canonical example: a coalition of secular
	leftists, liberals, and Islamists overthrew the Shah, whereupon the Islamist
	component, which had the clearest and most operationally concrete vision of the
	world it wanted, rapidly eliminated its former allies. The partners who had a precise
	positive project prevailed over those who had only a shared negation.
	
	\section{Consequences of not having a coherent positive vision}
	
	The most important conclusion that follows from this analysis is that, for a
	substantial portion of the contemporary Western left, there is no coherent and
	verifiable positive vision of the world they want to construct. What exists is a
	clear enemy, a coalition organised around opposition to that enemy, and a regulative
	ideal, the egalitarian utopia, that is sufficiently abstract to generate no
	operational commitments and therefore no internal disagreements.
	
	Coalitions organised primarily around negation have a characteristic weakness: they
	are sustained by the existence of the enemy and tend to dissolve when the enemy is
	gone, because the divergent positive values of the coalition's components then
	become visible and irreconcilable. In the case under analysis, the positive values
	of secular feminism and those of religious fundamentalism are not merely different;
	they are incompatible. The alliance holds as long as both parties can direct their
	energy toward the common adversary.
	
	The deeper question, what the political landscape would look like in the aftermath
	of the egalitarian utopia's construction, is one that the framework does not
	address, and arguably cannot address without fracturing the coalition that is
	supposed to bring it about. The energy of the movement is overwhelmingly diagnostic
	and critical, not constructive. It is more skilled at identifying and denouncing
	hierarchies than at specifying what non-hierarchical arrangements would look like or
	how they would be sustained.
	
	This is not, finally, a marginal or accidental feature. It is the structural
	consequence of building a political identity around guilt and opposition rather than
	around a positive account of justice. The world they want, to the extent that it
	exists at all, is most precisely described as the world that is left after the
	world they oppose has been dismantled. Whether what is left would be hospitable to
	the values, equality, freedom, secularism, in whose name the dismantling was
	undertaken is a question the framework is not equipped to answer.
	
	\section{Conclusion}
	
	The apparent incoherence of the contemporary Western left is not, this article has
	argued, the product of simple hypocrisy or collective irrationality. It is the
	product of a coherent internal logic that rests on three pillars: the replacement of
	class analysis with an identity-based power analysis that determines coalition
	membership by structural position rather than by shared values; a doctrine of
	collective guilt that frames Western civilisation as a permanent structure of
	injustice and renders the application of universal normative standards across
	cultural boundaries illegitimate; and a psychosocial dynamic in which the
	acknowledgement of that guilt functions as the constitutive identity of the
	political community, making it structurally interminable.
	
	Together these features produce a movement with a clear enemy, a broad if internally
	contradictory coalition, and no robust positive vision of the social order it would
	construct. The most serious risk this configuration generates is not to the
	adversaries the movement names. It is the risk, well-documented in the historical
	record, of being consumed by partners in the coalition who have a clearer and more
	operationally concrete answer to the question of what world they want.
	
\end{document}
